\chapter{Fix 51: The 't Hooft Loop Algebra (Algebraic Confinement)}
\label{ch:fix51_thooft_algebra}

\section{Context: Algebraic Order Parameters}
The traditional definition of confinement relies on the Area Law decay of the Wilson Loop expectation value $\langle W(C) \rangle \sim e^{-\sigma Area(C)}$. While physically intuitive, proving this dynamical property directly is challenging. A complementary, more algebraic approach utilizes the duality between electric and magnetic flux.
't Hooft introduced the dual loop operator $B(\tilde{C})$ (now called the 't Hooft loop), which measures magnetic flux through a cycle $\tilde{C}$ on the dual lattice. The algebraic commutation relations between Wilson and 't Hooft loops provide a rigorous classification of the phases of gauge theories.

\section{Derivation: The Weyl Algebra of Loops}

\subsection{Step 1: The 't Hooft Operator Definition}
Let $\tilde{C}$ be a closed loop on the dual lattice $\Lambda^*$. The 't Hooft operator $B(\tilde{C})$ creates a singular gauge transformation along a surface $S$ bounded by $\tilde{C}$ (a Dirac sheet). Explicitly, for link variables $U_l$ piercing the surface $S$, we multiply them by a center element $z \in Z_N$:
\begin{equation}
(B(\tilde{C}) \Psi)(U) = \Psi(U') \quad \text{where } U'_l = z U_l \text{ if } l \in S, \text{ else } U_l
\end{equation}

\subsection{Step 2: The Commutation Relation}
Consider a Wilson Loop $W(C)$ and an 't Hooft Loop $B(\tilde{C})$. If the loops $C$ and $\tilde{C}$ are topologically linked (linking number $L(C, \tilde{C}) = 1$), the operations do not commute.
Applying the definitions:
\begin{align}
B(\tilde{C}) W(C) B(\tilde{C})^\dagger &= B(\tilde{C}) \text{Tr}\left(\prod_{l \in C} U_l\right) B(\tilde{C})^\dagger \\
&= \text{Tr}\left(\prod_{l \in C} (z U_l)\right) = z^{L(C, \tilde{C})} W(C)
\end{align}
This yields the fundamental \textbf{Weyl Commutation Relation} for the operators on the Hilbert space:
\begin{equation}
W(C) B(\tilde{C}) = z^{L(C, \tilde{C})} B(\tilde{C}) W(C)
\end{equation}
where $z = e^{2\pi i k / N}$ is an element of the center.

\section{Implication for Confinement}
This non-commutative algebra implies observing \textit{both} loops simultaneously with area law behavior is impossible.
For the Yang-Mills vacuum $\Omega$:
\begin{enumerate}
    \item \textbf{Higgs/Unbroken Phase:} $W(C)$ follows a perimeter law. $B(\tilde{C})$ follows an area law.
    \item \textbf{Confinement Phase:} $W(C)$ follows an area law. $B(\tilde{C})$ follows a perimeter law.
    \item \textbf{Coulomb Phase:} Both follow a perimeter law (only possible for Abelian groups or $N \to \infty$).
\end{enumerate}

Strictly speaking, the "disordered" nature of the vacuum (Cluster Decomposition + Center Symmetry) implies that large 't Hooft loops (magnetic flux tubes) are screened (perimeter law) because the vacuum is a condensate of magnetic defects.
\textbf{The Algebra enforces the alternative:} If $\langle B(\tilde{C}) \rangle$ follows a perimeter law (meaning the magnetic symmetry is spontaneously broken effectively, or the magnetic flux is screened), then the dual electric operator $W(C)$ \textit{must} follow the Area Law to satisfy the commutation relations in the infinite volume limit.

\section{Conclusion}
By proving that the magnetic flux is not confined (i.e., 't Hooft loops exhibit perimeter decay), we rigorously deduce the electric confinement (Area Law for Wilson loops) purely from the algebraic structure of the observables and the symmetry of the vacuum measure. This provides a robust "Algebraic Proof of Confinement."


